\section{完整设计研究}
\label{app:design}

本附录的各项比较只改变一个轴，其他轴保持选定配置，采用附录~\ref{app:recipe} 的配方，各变体使用相同场景与推理种子。所有变体先用一个训练种子筛选，全部选择决策基于开发池，表中报告留出场景。编码器和耦合比较，以及同时出现在表~\mainref{tab:design} 中的行，平均三个训练种子；其余行使用一个种子。整个研究共耗费约 \fillin{XXXX} GPU 小时训练和 \fillin{XXXXXX} 条评测 episode。

\paragraph{选择。}
候选的部署性能须满足：从观测到动作的 p95 时间不超过 200\,ms，峰值显存低于单块 \cfgGPU{} 的 80\%。由于单种子结果可相差数个百分点（附录~\ref{app:statistics}），仅当候选在开发场景的种子 0 上提高超过 3 个百分点，且种子 1、2 均保持增益时，才替换参考配置；否则保留更便宜或更简单的变体。

\subsection{共用组件}
\label{app:design_shared}
观测编码器和仅训练时使用的外观目标由三项研究共用，我们先确定它们。开发场景上比较与策略共同训练的 CNN、冻结 DINOv3 ViT-S/16 与 ViT-L/16 \cite{dinov3}，以及 MAE ViT-B/16 \cite{mae}（表~\ref{tab:supp_shared}）。冻结 MAE 的三种子均值最高，为 \fillin{66.8}\%，但种子范围为 \fillin{63.7}--\fillin{71.1}\%，与 DINOv3 编码器重叠。我们依据其均值及低于 ViT-L/16 的成本选择它。DINOv3 ViT-L/16 特征仍作为未来外观目标之一，与世界表示一同选择（附录~\ref{app:design_predict}）。

% Transcribed from Research docs/experiments/SCALE0/SCALE0.md and AUX0/results/comparison.md (8-GPU encoder runs), BACKBONE0/BACKBONE0.md (matched parameter sharing). SR mean/sample SD from three seed counts /1400.
\begin{table}[!htb]
\centering\footnotesize
\caption{\tablelang{\textbf{Encoder and parameter sharing.} Three seeds, 1{,}400 scenes per seed; direct 3D, SE(3), and all appearance targets. SR: seed mean $\pm$ sample s.d.\ (\%). Panel (b) fixes MAE-B, reusing (a)'s shared-attention runs. Params: trainable backbone; $\star$: selected. Blue/bold and light blue mark the best and second point estimates within each panel and outcome column, not statistical significance.}{\textbf{编码器与参数共享。}三个种子，各评测 1{,}400 场；直接三维、SE(3) 与全部外观目标。SR 为种子均值 $\pm$ 样本标准差（\%）。面板 (b) 固定 MAE-B，复用 (a) 的共享注意力结果。参数量仅计可训练主干；$\star$ 表示选定配置。 各面板和结果列中，粗体加蓝底与浅蓝底分别标出最佳、次佳点估计，不表示统计显著性。}}
\label{tab:supp_shared}
\begin{tabularx}{\columnwidth}{@{}Xlc@{}}
\toprule
\tablelang{Configuration}{配置} & \tablelang{Status / params}{状态／参数量} & SR (\%) \\
\midrule
\multicolumn{3}{@{}l}{\tablelang{\textit{(a) Observation encoder}}{\textit{(a) 观测编码器}}}\\
CNN & \tablelang{trainable}{可训练} & \secondcell{$62.86\pm2.64$} \\
DINOv3 ViT-S/16 & \tablelang{frozen}{冻结} & $61.69\pm2.00$ \\
DINOv3 ViT-L/16 & \tablelang{frozen}{冻结} & $62.21\pm1.82$ \\
MAE ViT-B/16\sel & \tablelang{frozen}{冻结} & \bestcell{$66.76\pm3.84$} \\
\midrule
\multicolumn{3}{@{}l}{\tablelang{\textit{(b) World/action parameter sharing}}{\textit{(b) 世界／动作参数共享}}}\\
\tablelang{Shared attention\sel}{共享注意力\sel} & 47.266M & \bestcell{$66.76\pm3.84$} \\
\tablelang{Separate parameters}{独立参数} & 47.272M & \secondcell{$62.86\pm2.98$} \\
\bottomrule
\end{tabularx}
\end{table}


\subsection{记住}
\label{app:design_remember}
\paragraph{条件。}
各条件都接收当前图像、因果跟踪器给出的\emph{当前}目标位姿、八个机器人状态，以及选定的未来目标和耦合。差别仅在于保留过去 350\,ms 的什么信息，从而固定 $e_{\rm loc}$ 并隔离历史对 $e_v$ 的贡献。RGB 窗口保留每隔 50\,ms 的七帧历史，每帧 196 个图像块，不压缩。光流为 $64\times64$ RGB 上的双向 Farneb\"ack 流 \cite{farneback}，覆盖 100、100、100、50\,ms 四个区间，总计同一 350\,ms 窗口，带前向—后向有效性；“光流 + 运动统计”另加 DynamicWAM 每区间十二项物理统计 \cite{dynamicwam}。编码器运动向量来自实际 MPEG-4 P 帧，并按光流方式编码；视频隐变量为冻结 Wan2.2 VAE 的后验均值，由 56 个可学习查询读取。压缩窗口用共享的八查询组读取历史帧，或均匀分配，或按 PACT-WAM \cite{pactwam} 为较新帧分配更多查询。RGB-D 窗口增加每个历史图像块的度量位置，检验缺少的是度量几何还是运动。SE(3) 变体中，仅位置历史去掉朝向；检测器历史以逐帧单目三维框替代物体模型（附录~\ref{app:tracker}）；真值历史使用仿真器位姿。最后一个对照从 RGB 窗口中去掉当前目标位姿，测量当前位姿本身的贡献。

% Transcribed from Research HISTORY1/HISTORY1.md, completed nine-condition seed-0 study; original 4-layer synchronous evaluation.
\begin{table}[t]
\centering\footnotesize
\caption{\tablelang{\textbf{History screening.} Four layers, width 768, frozen DINO-L, one training seed, and the original synchronous evaluation on 1,400 scenes. The first eight conditions receive the current target pose; the final condition adds target-motion history. These are not final-scale real-time results.}{\textbf{历史筛选。}4层、宽768、冻结DINO-L、一个训练种子，按原同步协议评测1,400个场景。前八项接收当前目标位姿，最后一项另加目标运动历史。这些不是最终规模的实时结果。}}
\label{tab:supp_remember}
\begin{tabular}{@{}lrr@{}}
\toprule
\tablelang{History input}{历史输入} & \tablelang{Success / 1,400}{成功数 / 1,400} & SR (\%) \\
\midrule
\tablelang{Current image}{当前图像} & 748 & 53.43 \\
\tablelang{RGB window}{RGB窗口} & 662 & 47.29 \\
\tablelang{Optical flow}{光流} & 720 & 51.43 \\
\tablelang{Flow + motion statistics}{光流＋运动统计} & 723 & 51.64 \\
\tablelang{Codec motion vectors}{编码器运动向量} & 733 & 52.36 \\
\tablelang{Video-VAE latent}{视频VAE隐变量} & 710 & 50.71 \\
\tablelang{Compressed window (age)}{压缩窗口（按时间）} & 634 & 45.29 \\
\tablelang{Compressed window (uniform)}{压缩窗口（均匀）} & 670 & 47.86 \\
\midrule
\tablelang{Current image + target-motion history}{当前图像＋目标运动历史} & 824 & 58.86 \\
\bottomrule
\end{tabular}
\end{table}


\paragraph{冻结读出。}
为区分历史编码\emph{包含}的信息与策略\emph{使用}的信息，我们在三处冻结特征上拟合岭回归读出：历史编码器输出、最终物体查询、首次函数求值时动作头的输入。标签为过去 50、200、350\,ms 的目标位置、最后区间的速度及 400\,ms 未来位移，均在基座坐标系中。每项任务分别用 60、20、20 条不同开发 episode 训练、验证和测试，每条五个窗口。依据验证误差从 $\{0.1,1,10,100\}$ 选择正则强度，以打乱标签的对照估计随机水平。条件读出先从当前图像、几何、机器人状态与当前位姿拟合基线，再从相应位置的特征预测其残差，以检验历史超出现有状态的增量信息。

\paragraph{扰动与干预。}
离线扰动将所有历史槽位替换为当前图像，或倒序历史图像，同时固定当前帧、位姿和噪声，测量首条发出动作的变化。在最终模型上（表~\ref{tab:supp_interventions}），\emph{静止历史}在重算运动学前将全部历史位姿替换为当前位姿，消除位移、速度和加速度。\emph{时间反转}倒序历史位姿。\emph{当前状态匹配替换}使用同任务、物体和相机的另一条最近开发 episode 的历史，要求当前目标和末端位姿在 5\,mm、$5^\circ$ 内。\emph{方向歧义}子集进一步要求末区间速度方向相反，且大小至少为 0.01\,m/s。

% Transcribed from Research HISTORY1/HISTORY1.md, frozen GT-E intervention (2026-10-18), original 4-layer seed-0 synchronous evaluation.
\begin{table}[t]
\centering\footnotesize
\caption{\tablelang{\textbf{Stationary-history intervention.} One frozen 4-layer HISTORY1 checkpoint, training seed 0, and its original synchronous protocol. Past target poses are replaced by the current pose and kinematics recomputed; the unchanged-history result is reused. These 1,400 original scenes are separate from the final real-time confirmation.}{\textbf{静止历史干预。}使用同一个冻结的4层 HISTORY1 检查点、训练种子0及其原同步协议。将过去目标位姿替换为当前位姿，再重算运动学；真实历史结果直接复用。这1,400个原场景与最终实时确认实验分开报告。}}
\label{tab:supp_interventions}
\begin{tabular}{@{}lrr@{}}
\toprule
\tablelang{Target history}{目标历史} & \tablelang{Success / 1,400}{成功数 / 1,400} & SR (\%) \\
\midrule
\tablelang{Unchanged}{真实历史} & 824 & 58.86 \\
\tablelang{Stationary replacement}{静止历史替换} & 9 & 0.64 \\
\bottomrule
\end{tabular}
\end{table}


\paragraph{位姿噪声。}
计算运动学前，对各历史位姿独立加入位置标准差 $\sigma_p\in\{0,1,2,5,10,20,40\}$\,mm、旋转标准差 $\sigma_R=\sigma_p/(30\,\mathrm{mm})$ rad 的高斯噪声，使两者在 3\,cm 力臂处产生同样位移。当前位姿保持不变，从而隔离历史作用。取 $T=350$\,ms、$n=8$ 及实际领先时长，依据表~\ref{tab:roles} 逐任务计算预测交叉点 $\sigma^\star$（附录~\ref{app:predictions}），并与跟踪器实测误差（表~\ref{tab:supp_tracker}）一起标在图~\mainref{fig:analysis}c。

\paragraph{遮挡与采样。}
用 0.5、1、2、3\,s 遮挡及无遮挡对照检验历史缺口。采样试验保持 350\,ms 跨度，选择 2、4 或 8 个偏移，时间嵌入使用实际偏移。

\subsection{预测}
\label{app:design_predict}
\paragraph{表示。}
表~\ref{tab:target_shapes} 定义各类未来目标。所有目标均相对当前状态累积，覆盖从 50 到 400\,ms 的相同八个时域。表~\ref{tab:supp_predict} 各行保持选定配置的历史、耦合、动作损失、训练预算和世界流容量，仅替换由场景查询经附录~\ref{app:world} 接口回归的世界目标，因此图像空间目标在容量匹配下比较。冻结视频一行则读取预训练图像到视频模型从当前帧预测的隐变量，其未来不依赖我们的模型容量。Kabsch 行先预测物体点轨迹，再拟合 SE(3)，所获信息与点轨迹行相同，但增加刚性先验。最后一组保留选定世界，仅改变部署时删除读出头的外观辅助目标，以测量它们教给场景查询的内容。

开发场景中，多数表示的总体差异只有数个百分点，差别集中于少数任务。目标点轨迹配合流的总体 SR 甚至略高于 SE(3) 轨迹（单种子 \fillin{+1.7} 个百分点）。我们保留直接给出位姿、且误差预算在旋转目标上支持的 SE(3) 轨迹，表~\ref{tab:supp_predict} 在留出场景上比较两者。稠密二维位移是唯一明显有害的未来目标（相对无未来为 \fillin{$-$3.6} 个百分点）。

\begin{table*}[t]
\centering\small
\caption{\tablelang{Future targets in the representation screen. Shape excludes batch/query axes.
Displacements are cumulative from the current state except codec vectors,
which retain the decoder's native correspondence. All targets come from
the recorded demonstration.}{表示筛选中的未来监督目标。张量形状不含 batch 和 query 轴。除保留解码器原生对应关系的 codec 向量外，位移均从当前状态累计。所有监督目标均来自记录的演示。}}
\label{tab:target_shapes}
\begin{tabularx}{\textwidth}{@{}llX@{}}
\toprule
\tablelang{Target}{目标} & \tablelang{Shape}{形状} & \tablelang{Support, units, and validity}{几何支持、单位与有效条件} \\
\midrule
\tablelang{Scene flow}{场景流} & $196\times8\times3$ & \tablelang{Fixed $14\times14$ patch centers; base-frame meters; valid source depth/geometry and finite future displacement.}{固定 $14\times14$ patch 中心；以基座坐标系中的米为单位；源深度／几何有效，未来位移有限。} \\
\tablelang{Target pose}{目标位姿} & $8\times6$ & \tablelang{Cumulative translation (m) and left-relative rotvec (rad); valid object pose and query.}{累计平移 (m) 和左相对旋转向量 (rad)；物体位姿和 query 有效。} \\
P04: Wan & $48\times2\times12\times10$ & \tablelang{Frozen Wan2.2 VAE posterior mean and native normalization; future groups end at 200/400\,ms; valid query.}{冻结的 Wan2.2 VAE 后验均值及原生归一化；未来分组终点为 200/400\,ms；query 有效。} \\
\tablelang{P05: dense 2-D}{P05：稠密二维} & $8\times2\times224\times224$ & \tablelang{Fixed source pixels; displacement divided by 224; valid source and finite future projection.}{固定源像素；位移除以 224；源像素有效，未来投影有限。} \\
\tablelang{P06: codec}{P06：codec} & $8\times2\times64\times64$ & \tablelang{Each future/current pair is encoded as P/I frames; native vectors divided by 64; actual valid vector regions.}{每对未来／当前帧编码为 P/I 帧；原生向量除以 64；使用实际有效向量区域。} \\
\tablelang{P07: scene tracks}{P07：场景轨迹} & $196\times8\times3$ & \tablelang{Same patch source points; two displacement channels divided by 224 and one visibility logit; valid correspondences.}{相同的 patch 源点；两个位移通道除以 224，另加一个可见性 logit；对应关系有效。} \\
\tablelang{P08: object DINO}{P08：物体 DINO} & $8\times1024$ & \tablelang{Object-mask-area-weighted patch pooling; channel layer normalization; positive mask area.}{按物体掩码面积加权汇聚 patch；沿通道作层归一化；掩码面积为正。} \\
\tablelang{P09: object tracks}{P09：物体轨迹} & $32\times8\times3$ & \tablelang{Deterministic farthest-point sampling in the current target mask; displacement and visibility as P07; padded points invalid.}{在当前目标掩码内确定性地作最远点采样；位移和可见性同 P07；填充点无效。} \\
\tablelang{$A_\alpha$: magnitude}{$A_\alpha$：幅值} & $196\times8$ & \tablelang{Cumulative flow norm (m); same correspondence mask as flow.}{累计流的范数 (m)；对应关系掩码与场景流相同。} \\
\tablelang{P12: depth}{P12：深度} & $8\times1\times224\times224$ & \tablelang{Metric future depth; finite positive pixels.}{未来深度的米制值；像素值为有限正数。} \\
P13: DINO & $8\times196\times1024$ & \tablelang{Frozen DINOv3-L/16 patches, channel layer normalization; valid query/horizon/patch.}{冻结的 DINOv3-L/16 patch，沿通道作层归一化；query、时域和 patch 均有效。} \\
\tablelang{P14: dynamic RGB}{P14：动态 RGB} & $8\times3\times224\times224$ & \tablelang{RGB divided by 255; visible surfaces moving more than 1\,mm in world coordinates, including robot and other geometry.}{RGB 除以 255；世界坐标中移动超过 1\,mm 的可见表面，包括机器人和其他几何体。} \\
\bottomrule
\end{tabularx}
\end{table*}

% Transcribed from tables/future_complete.tex (Research FUTURE0 accepted summaries) and Research AUX0/AUX0.md MAE-B seed-1 results (2026-10-22).
\begin{table*}[p]
\centering\small
\caption{\tablelang{\textbf{Forecast targets} at the screening scale. The first group uses a frozen DINO-L encoder and seed 0, with its 140 development scenes and the 1{,}260 remaining scenes reported separately. The second group uses the MAE encoder and seed 1 on 1{,}400 scenes; it keeps scene flow, SE(3), and the current-target RGB reconstruction and removes future appearance targets. Compare conditions within each group.}{\textbf{预测目标}（筛选规模）。第一组使用冻结 DINO-L 编码器与种子 0，其 140 个开发场景与其余 1{,}260 个场景分别报告。第二组使用 MAE 编码器与种子 1，在 1{,}400 个场景上评测；它保留场景流、SE(3) 与当前目标 RGB 重建，并移除未来外观目标。请在组内比较各条件。}}
\label{tab:supp_predict}
\setlength{\tabcolsep}{4pt}
\begin{tabular*}{\textwidth}{@{\extracolsep{\fill}}llrrr@{}}
\toprule
\tablelang{ID}{编号} & \tablelang{Future supervision}{未来监督} & \tablelang{Dev.\ success / 140}{开发集成功数 / 140} & \tablelang{Success / 1,260}{成功数 / 1,260} & SR$_{1260}$ (\%) \\
\midrule
P00 & \tablelang{None}{无} & 88 & 791 & 62.78 \\
P01 & \tablelang{3D scene flow}{三维场景流} & 95 & 782 & 62.06 \\
P02 & \tablelang{SE(3)}{SE(3)} & 92 & 789 & 62.62 \\
P03 & \tablelang{Scene flow + SE(3)}{场景流＋SE(3)} & 93 & 786 & 62.38 \\
P04 & \tablelang{Video-VAE latent + SE(3)}{视频VAE隐变量＋SE(3)} & 92 & 783 & 62.14 \\
P05 & \tablelang{Dense 2D displacement + SE(3)}{稠密二维位移＋SE(3)} & 87 & 746 & 59.21 \\
P06 & \tablelang{Codec motion + SE(3)}{编码器运动向量＋SE(3)} & 86 & 795 & 63.10 \\
P07 & \tablelang{Sparse scene 2D tracks + SE(3)}{稀疏场景二维轨迹＋SE(3)} & 90 & 807 & 64.05 \\
P08 & \tablelang{Scene flow + object DINO}{场景流＋物体DINO} & 89 & 794 & 63.02 \\
P09 & \tablelang{Scene flow + target 2D tracks}{场景流＋目标二维轨迹} & 88 & 808 & 64.13 \\
$A_\alpha$ & \tablelang{P03 + motion magnitude}{P03＋运动幅值} & 89 & 778 & 61.75 \\
P12 & \tablelang{P03 + depth}{P03＋深度} & 90 & 791 & 62.78 \\
P13 & \tablelang{P03 + DINO}{P03＋DINO} & 92 & 788 & 62.54 \\
P14 & \tablelang{P03 + dynamic RGB}{P03＋动态RGB} & 89 & 790 & 62.70 \\
P15 & \tablelang{P03 + depth / DINO / dynamic RGB}{P03＋深度／DINO／动态RGB} & 94 & 802 & 63.65 \\
\midrule
\multicolumn{2}{@{}l}{\tablelang{AUX0 MAE-B: future appearance auxiliaries}{AUX0 MAE-B：未来外观辅助}} & \multicolumn{2}{c}{\tablelang{Success / 1,400}{成功数 / 1,400}} & SR$_{1400}$ (\%) \\
\midrule
\multicolumn{2}{@{}l}{\tablelang{None}{无}} & \multicolumn{2}{c}{894} & 63.86 \\
\multicolumn{2}{@{}l}{\tablelang{Depth only}{仅深度}} & \multicolumn{2}{c}{867} & 61.93 \\
\multicolumn{2}{@{}l}{\tablelang{DINO only}{仅DINO}} & \multicolumn{2}{c}{946} & 67.57 \\
\multicolumn{2}{@{}l}{\tablelang{Dynamic RGB only}{仅动态RGB}} & \multicolumn{2}{c}{926} & 66.14 \\
\multicolumn{2}{@{}l}{\tablelang{All three (existing reference)}{三者全部（已有参照）}} & \multicolumn{2}{c}{995} & 71.07 \\
\bottomrule
\end{tabular*}
\end{table*}


\paragraph{角色覆盖。}
在选定表示上，各角色行保留全部源点和相同计算量，仅改变哪些角色参与流损失。最后一行去掉角色嵌入，检验显式角色划分是否在依赖任务的读出之外提供帮助。

\paragraph{分辨率与间隙。}
$28^2$、$56^2$ 网格更密集地采样同一图像；按物体采样的流将点放在环境掩码上，而非固定网格。所有变体使用相同标签、归一化与损失权重。路径误差 $e_F$ 是最接近执行路径的 256 个留出环境点，在路径经过它们时的流端点误差。表~\ref{tab:supp_clearance} 报告全部变体在附录~\ref{app:sweeps} 间隙扫描上的结果，并加入容量匹配的视频隐变量未来。

% Planned results: context resolution vs clearance on the passage tasks (App. D.3; Fig. 4e data), held-out scenes, three seeds (Fig. 4e).
\begin{table}[t]
\centering\footnotesize
\caption{\tablelang{\textbf{Context resolution versus clearance} on the passage tasks: grid spacing at the workspace and path
error $e_F$ (mm), and SR (\%) at the swept clearances $c_1<\dots<c_5$ (Appendix~\ref{app:sweeps}). The target is predicted in
SE(3) throughout.}{\textbf{环境分辨率与间隙}（穿越任务）：工作空间中的网格间距与路径误差 $e_F$（mm），以及各扫描间隙 $c_1<\dots<c_5$
下的成功率（\%）（附录~\ref{app:sweeps}）；目标始终以 SE(3) 预测。}}
\label{tab:supp_clearance}
\setlength{\tabcolsep}{2.3pt}
\begin{tabular}{@{}lccccccc@{}}
\toprule
\tablelang{Context}{环境} & \tablelang{Spacing}{间距} & $e_F$ & $c_1$ & $c_2$ & $c_3$ & $c_4$ & $c_5$ \\
\midrule
\tablelang{none}{无} & -- & -- & \resultplaceholder & \resultplaceholder & \resultplaceholder & \resultplaceholder & \resultplaceholder \\
\tablelang{flow, $14^2$}{流，$14^2$} & \resultplaceholder & \resultplaceholder & \resultplaceholder & \resultplaceholder & \resultplaceholder & \resultplaceholder & \resultplaceholder \\
\tablelang{flow, $28^2$}{流，$28^2$} & \resultplaceholder & \resultplaceholder & \resultplaceholder & \resultplaceholder & \resultplaceholder & \resultplaceholder & \resultplaceholder \\
\tablelang{flow, $56^2$}{流，$56^2$} & \resultplaceholder & \resultplaceholder & \resultplaceholder & \resultplaceholder & \resultplaceholder & \resultplaceholder & \resultplaceholder \\
\tablelang{flow, object-sampled}{流，按物体采样} & \resultplaceholder & \resultplaceholder & \resultplaceholder & \resultplaceholder & \resultplaceholder & \resultplaceholder & \resultplaceholder \\
\tablelang{video-VAE latent}{视频 VAE 隐变量} & -- & -- & \resultplaceholder & \resultplaceholder & \resultplaceholder & \resultplaceholder & \resultplaceholder \\
\bottomrule
\end{tabular}
\end{table}


\paragraph{外推基线。}
匀速和匀加速预测只使用模型可见的因果 SE(3) 历史，即末区间速度与最后两个区间的速度差，分别作用于平移和以旋转向量表示的旋转。离线时，将三类预测与生效时刻真值比较。为测量闭环价值，以相同数据和预算训练两个变体（表~\ref{tab:supp_extrapolation}）：第一个把外推目标轨迹作为额外输入 token，保留选定模型的世界监督；第二个以这些输入 token 替代世界监督，即表~\mainref{tab:main} 的 CV-input 基线。

% Planned results: offline causal CV/CA and frozen learned-pose-head errors; no new controller or training arm.
\begin{table}[!htb]
\centering\footnotesize
\caption{\tablelang{\textbf{Offline prediction errors.} Target-position error at 400\,ms (mm) on the same held-out observations. CV and CA are analytic extrapolations from causal pose history; the learned head reads frozen full-model states.}{\textbf{离线预测误差。}在同一批留出观测上计算400\,ms处的目标位置误差（mm）。CV和CA是由因果位姿历史计算的解析外推；学习位姿头读取冻结完整模型的状态。}}
\label{tab:supp_extrapolation}
\setlength{\tabcolsep}{2.6pt}
\begin{tabular}{@{}lcccc@{}}
\toprule
 & \tablelang{Const.\ vel.}{匀速} & \tablelang{Harmonic}{谐波} & \tablelang{Roll/slide}{滚动/滑动} & \tablelang{Impact}{撞击} \\
\midrule
\tablelang{CV extrapolation}{匀速外推} & \resultplaceholder & \resultplaceholder & \resultplaceholder & \resultplaceholder \\
\tablelang{CA extrapolation}{匀加速外推} & \resultplaceholder & \resultplaceholder & \resultplaceholder & \resultplaceholder \\
\tablelang{Learned pose head}{学习的位姿头} & \resultplaceholder & \resultplaceholder & \resultplaceholder & \resultplaceholder \\
\bottomrule
\end{tabular}
\end{table}


\paragraph{预测误差。}
物理误差使用共同的 256 个目标源点和 256 个环境源点，以确定性方式采样并排除在训练查询之外。平移与旋转误差分别保留度量单位和测地单位，流误差为有效对应上的端点误差。所有预测误差均在留出 episode 的生效时刻测量。

\subsection{行动}
\label{app:design_act}
\paragraph{主干。}
共享注意力主干在所有 token 组间共享注意力投影，按组分开前馈网络。独立参数组主干的世界与动作组各自包含注意力投影、自适应层归一化和前馈网络；条件与世界 token 共用世界组。通过将独立组前馈扩展比由 4 改为 3 匹配可训练参数量，两者均约 47.3M，层数、宽度和注意力头数相同。相同训练预算不意味着相同 FLOPs，因此报告实测成本。

\paragraph{耦合。}
各耦合的注意力掩码定义见附录~\ref{app:coupling}，模块、初始化、损失和训练预算相同。R0 保留世界查询与读出头，但关闭全部世界损失，从而与 R1 容量匹配且无世界监督。R1 加入世界监督，关闭所有世界到动作的边。其同一检查点分别仅运行动作，或计算但不读取预测，以隔离无用预测的成本。R2 让动作读取世界查询隐藏状态，R3 让世界读取带噪动作，R4 打开双向连接。R3、R4 每个 Euler 步重算世界。R5 解码世界显式几何，再通过 0.37M 参数的几何编码器重新编码供动作使用，训练端到端穿过预测几何，不采用教师强制。由于 R5 同时改变动作读取的内容及其进入位置，R2 与 R5 比较的是完整设计，而非单因素。

\Qthree{3}{开发场景中，没有耦合稳定优于 R1。世界读取动作的 R3 三种子均值最高，比 R1 高 \fillin{+0.6} 个百分点，但差异完全来自线缆通行，且两个读取方向的作用均跨种子变号。我们保留了同时也是最便宜的 R1。世界监督本身，即 R1 相对 R0，平均提高 \fillin{3.1} 个百分点，但各种子增益从 \fillin{$-$1.3} 到 \fillin{+9.6} 个百分点。}

% Planned results: Q3-INFO0 4x768, MAE, 4 epochs, three training seeds per R0-R5. Screening layout; no final-scale R5 run or equivalence claim.
\begin{table}[t]
\centering\footnotesize
\caption{\tablelang{\textbf{Couplings R0--R5 and backbones} at the screening scale (MAE encoder, three training seeds each; Appendix~\ref{app:design_act}).}{\textbf{耦合方式 R0--R5 与主干}（筛选规模：MAE 编码器，各三个训练种子；附录~\ref{app:design_act}）。}}
\label{tab:supp_act}
\setlength{\tabcolsep}{4.4pt}
\begin{tabular}{@{}lccc@{}}
\toprule
\tablelang{Coupling}{耦合} & \tablelang{SR mean}{SR 均值} & \tablelang{SR std.}{SR 标准差} & p95 (ms) \\
\midrule
R0 \tablelang{(future losses off)}{（关闭未来损失）} & \resultplaceholder & \resultplaceholder & \resultplaceholder \\
R1 \tablelang{(queries, no edge)}{（查询，无连接）} & \resultplaceholder & \resultplaceholder & \resultplaceholder \\
R2 \tablelang{(world$\to$action)}{（世界$\to$动作）} & \resultplaceholder & \resultplaceholder & \resultplaceholder \\
R3 \tablelang{(action$\to$world)}{（动作$\to$世界）} & \resultplaceholder & \resultplaceholder & \resultplaceholder \\
R4 \tablelang{(bidirectional)}{（双向）} & \resultplaceholder & \resultplaceholder & \resultplaceholder \\
R5 \tablelang{(decoded geometry)}{（解码几何）} & \resultplaceholder & \resultplaceholder & \resultplaceholder \\
\bottomrule
\end{tabular}
\end{table}


\paragraph{预测干预。}
干预固定条件 token 和动作噪声，仅改变动作从世界读取的内容。对 R2，替换动作 token 在每层读取的键和值；当前状态匹配替换使用当前目标位姿在 5\,mm 和 $5^\circ$ 内的另一条开发 episode 的世界状态，时间平移则将所有时域偏移 50\,ms 后重算。对 R5，在几何编码前编辑解码几何：零运动采用零三维位移与单位目标位姿，保持可见性不变；匹配当前状态的打乱采用同样匹配条件下另一条开发 episode 的几何；时间平移在 $\delta_h\pm50$\,ms 重采样轨迹；真值干预使用首次接触前的实测外源运动。接触后未来标签依赖动作，因此真值分析从仿真器快照分叉，在固定候选动作序列下进行。恒等对照复制原始状态，动作完全不读取世界的耦合充当阴性对照。
